\section{Analysis of the member paths}\label{sec:analysis}
The benchmark measures the complete GNNM recipe, which changes both the projectors and the sharing of graph weights. To examine the contribution of sharing itself, we start with GNNM and give each member its own copy of the graph stack. The copies initially contain the same values, so the member functions are matched before training. During training, the copies can change independently. Both versions choose a checkpoint using the same averaged validation prediction. Their difference measures what happens when the requirement to share graph weights is removed, under this training recipe.

We make this comparison on a favorable benchmark setting: five-layer SAGE, width 512, on Roman Empire masks 0--4. Every four-member method in Table~\ref{tab:component-main} uses the common checkpoint rule. The setting was chosen after the benchmark, so the study examines one observed success. It does not estimate the effect of sharing across new graphs.

\begin{table}[t]
\centering\small
\setlength{\tabcolsep}{5pt}
\caption{Roman Empire component study. Accuracy is mean $\pm$ standard deviation over five overlapping official masks. All four-member rows use a pooled checkpoint rule. The first two rows form the matched-initialization sharing comparison. The remaining component comparisons have the initialization relationships specified in Appendix~\ref{app:controls}.}
\label{tab:component-main}
\begin{tabular}{lrr}
\toprule
Method & Accuracy (\%) & Parameters ($10^6$)\\
\midrule
GNNM & $89.50\pm0.29$ & 6.738\\
Private propagation for every member & $86.56\pm0.57$ & 26.437\\
\midrule
Four ordinary models (pooled ENS) & $87.52\pm0.26$ & 26.923\\
Four independent input/output maps & $86.99\pm0.38$ & 7.221\\
Input factors only & $89.35\pm0.32$ & 6.736\\
Output factors only & $84.96\pm0.47$ & 6.733\\
Four output heads & $84.97\pm0.56$ & 6.758\\
Frozen output factors & $89.48\pm0.28$ & 6.736\\
Factors also inside graph layers & $89.66\pm0.14$ & 6.871\\
One-member GNNM & $85.47\pm0.64$ & 6.732\\
BASE & $86.21\pm0.47$ & 6.731\\
\bottomrule
\end{tabular}
\end{table}

GNNM exceeds the version with private graph stacks by 2.94 accuracy points on average, with positive differences on all five masks. It also exceeds the ordinary ensemble under the common checkpoint rule by 1.99 points. The private-stack model can reproduce any GNNM prediction by keeping its four graph stacks identical. GNNM's higher accuracy here therefore cannot come from access to a larger set of prediction functions. It concerns what this training and checkpoint-selection procedure finds under the sharing constraint. Giving members private weights changes both capacity and optimizer histories, which this comparison does not separate. Repeating the matched sharing comparison on mask 0 with three optimization seeds gives a gain of $3.33\pm0.63$ points (sample SD), positive for every seed. This remains evidence from one graph and a finite training budget.

Sharing weights is only part of GNNM: its members also enter the graph layers with different hidden states. To examine this distinction, we compare varying the input map with varying only the output map. The latter makes several predictions from one common propagated representation. The input-only variant scores close to GNNM, whereas output-only and heads-only variants score lower in this case. Freezing the output factors barely changes the mean, while adding factors inside every graph layer gives a small increase. This pattern is consistent with useful input variation in this setting. Several variants start from different functions, so these comparisons do not isolate a single component or establish that both boundary projectors are necessary or optimal.

A higher ensemble score can arise because its members become more accurate, because averaging corrects more of their errors, or both. We therefore return to the saved predictions from the matched sharing comparison. For each method, we calculate the mean accuracy of its four members and then the additional accuracy obtained by averaging their raw scores. This divides the observed ensemble difference into two quantities without changing any predictions.

For the Roman SAGE comparison, GNNM's members are 3.08 points more accurate than the untied members on average. Its gain from averaging is 0.14 points smaller. Together these changes give the 2.94-point pooled advantage. Pairwise class disagreement falls from 8.00\% to 6.07\%. Thus this favorable result comes from better members despite lower measured disagreement. More disagreement alone is not a sufficient explanation of ensemble quality.

A parallel comparison with Roman GAT checks whether this SAGE pattern also appears in an attention backbone. The mean shared-minus-private difference is only 0.04 points, with mixed signs across masks. The original GAT losses therefore do not reproduce as a large penalty from sharing in this case. Together, these results locate a favorable SAGE effect and prevent its member-accuracy explanation from being applied to every backbone. They describe the selected predictions, without establishing why training finds them. Appendix~\ref{app:gat-pair} gives the GAT masks and decision counts.

